\chapter{Significance of the Study}

This study matters to the three groups who live inside a learning platform, and to the people who build them.

\section{Benefits to Students}

Take the most obvious case: a student opening the course. Instead of a gradebook, they get the whole syllabus as a mindmap. What they have finished is marked completed, what they have open is in progress, what comes later is still locked. Reading that takes a glance. Watching a node change state after a resource is finished gives the kind of small, immediate feedback that game designers have relied on for decades; Hamari et al. found that game-like progress systems can raise motivation and engagement in educational settings \cite{hamari2014does}. The learning-path view adds a second reading of the same data, arranged as a path instead of a graph, for students who find the canvas noisy.

\section{Benefits to Teachers}

For a teacher the change is smaller to describe and larger to feel. A single page lists every student in the class, sorted into progressing, stuck or inactive by how recently they touched the material. No scrolling through rows to work out who is slipping. The same page places the class position against the syllabus position, so a topic the teacher has finished covering but the class has not started is visible as a gap instead of a hunch.

\section{Benefits to Educational Institutions}

One deployment can serve several schools or departments. Data from one institution is not readable from another, which is checked at the database instead of asserted in a policy document. An operator account sees tenants; everyone else sees their own. For a college running separate installations per department, that is a practical saving in both cost and maintenance.

\section{Contribution to the EdTech Field}

Three ideas usually studied apart are combined here: mindmap-based progress visualisation, dual progress tracking, and retrieval scoped to a course. Each has its own literature. Working systems that carry all three are rare, and the interesting question is whether they cooperate or interfere. Synapse is an attempt to answer that with something that runs.

\section{Broader Societal Impact}

Wider than any single course, better progress tools move intervention earlier. A student who is caught in week three is cheaper to help than one who is caught after the exam. The stack is open source throughout, so a department can read the code, change it, or run it themselves.
